\section{Polynomial encoding with few nonlinear inequalities}
\label{sec:fixedcount}
This section proves a second upper bound: if at most $k$ native
inequalities are nonlinear in the continuous variables, some sufficient
coefficient has $N^{O(k+1)}$ bits. For fixed $k$ this is polynomial in the
input length, however large the continuous dimension, the number of affine
rows or the number of integer assignments. In practical terms, many linear
rows together with a few convex quadratic rows cannot produce the doubly
exponential growth of Section~\ref{sec:lower}.

We work in the setting of Section~\ref{sec:upper}. Let $k$ be an upper
bound on the number of native inequalities $g_i(x,z)\le0$ that are
nonlinear in the continuous variables $x$. Constraints that are affine in
$x$ are not counted, even if they contain integer variables, such as
$x_1z_1+z_1^2\le3$; the objective is not counted either. As noted after
Assumption~\ref{ass:model}, whether a row is nonlinear in $x$ does not
depend on $z$.

\begin{theorem}\label{thm:fixedcount}
Suppose that Assumption~\ref{ass:model} holds and that at most $k$
native inequalities are nonlinear in the continuous variables. Then there is an integer $\rho\ge1$ with
\begin{equation}\label{eq:fixedbits}
 \operatorname{bit}(\rho)\le N^{O(k+1)}
\end{equation}
that is minimizer-exact at multiplier zero for both the infinity norm and
the one-norm; the same $\rho$ works for both. The constant in the exponent
is absolute.
\end{theorem}
Neither Theorem~\ref{thm:generalupper} nor Theorem~\ref{thm:fixedcount}
implies the other: the first is polynomial for fixed $n$, the second for
fixed $k$.

\paragraph{Proof strategy.}
In the proof of Theorem~\ref{thm:generalupper}, both the separation of
equality-infeasible slices and the linking multipliers of
equality-feasible slices are bounded through radii of semialgebraic sets
in about $2n$ variables, which costs a factor $2^{O(n)}$ in the exponent.
Here each slice needs only one scalar. For an equality-infeasible slice
it is the separation $\delta_z=\min_{X_z}\norm r_\infty$. For an
equality-feasible slice it is a Slater margin $\gamma_z$, the largest
amount (capped at one) by which some point of the slice with $r=0$
satisfies all nonlinear rows strictly. Both are optimal values of convex quadratic problems with
at most $k$ nonlinear rows. Lemma~\ref{lem:nonzerovalue} below bounds such
nonzero optimal values below by $2^{-N^{O(k+1)}}$; here $k$ enters only
the exponent of $N$, and $n$ does not enter it.
The Slater margin then bounds the multipliers of the nonlinear rows, and
the linking multipliers are recovered from the remaining affine
stationarity equations by linear algebra (Cramer's rule).

\subsection{An algebraic bound on optimal values}
The lemma says that the optimal value of a convex quadratic problem with
$k$ quadratic rows is an algebraic number whose degree and coefficient
size are polynomial in the input for fixed $k$; a nonzero such number
cannot be smaller than $2^{-\mathrm{poly}}$.

\begin{lemma}\label{lem:nonzerovalue}
Minimize a rational convex quadratic function over a nonempty set
described by a rational polyhedron with finite coordinate bounds and at
most $k$ rational convex quadratic inequalities. Let $N\ge2$ be the input
length of this problem, encoded explicitly as in
Assumption~\ref{ass:model}, and let $\omega$ be its optimal value. Then
$\omega$ is a root of a nonzero univariate integer polynomial whose degree
and coefficient bitsize are at most $N^{O(k+1)}$. In particular, if
$\omega\ne0$, then
\begin{equation}\label{eq:nonzero}
 |\omega|\ge2^{-N^{O(k+1)}}.
\end{equation}
No Slater condition or other constraint qualification is required.
\end{lemma}
\begin{proof}
The proof has three steps. (1) We delete the affine inequalities that are
inactive at an optimum and substitute the active affine face; this leaves
an unconstrained affine parametrization and at most $k+1$ quadratic rows.
(2) We regularize the objective and relax the rows by a parameter
$\eta>0$, so that KKT conditions hold with a positive definite matrix, and
eliminate the primal variables with the adjugate; this gives polynomial
conditions in $\eta$, the value $w$ and the multipliers. (3) We express
the limit $\eta\downarrow0$ by a formula with two quantifier blocks, of
sizes $1$ and at most $k+3$, and apply effective quantifier elimination,
whose degree and bitsize bounds are exponential only in the block sizes.

Write the problem in variables $x\in\R^n$ as the minimization of
$\phi_0(x)$ over the points of the polyhedron $P$ that satisfy
$\phi_i(x)\le0$ for $1\le i\le k$. The feasible set is nonempty and
compact, so an optimum $x^*$ exists. Note that $n\le N$.

\emph{Step 1: affine-face reduction.}
Let $E$ be the affine space defined by all equalities of $P$ and all
affine inequalities of $P$ that are active at $x^*$. The remaining affine
inequalities of $P$ have positive slack at $x^*$, hence on a neighborhood
$W$ of $x^*$. Deleting them leaves $x^*$ optimal over
$E\cap\{\phi_i\le0,\ 1\le i\le k\}$. Indeed, suppose a point $y$ of this
set had $\phi_0(y)<\phi_0(x^*)$. By convexity, every point
$x^*+\theta(y-x^*)$ with $0<\theta\le1$ lies in $E$, satisfies the
inequalities $\phi_i\le0$, and has objective value below $\phi_0(x^*)$.
For small $\theta$
it also lies in $W$, so it satisfies the deleted rows and is feasible for
the original problem. This contradicts the optimality of $x^*$.

Choose independent equations defining $E$. Gaussian elimination writes
$E=\{x_0+Vu:u\in\R^{n'}\}$, where $x_0$ and $V$ are rational with bit
length polynomial in $N$, and the free coordinates $u$ are $n'$ of the
coordinates of $x$; thus $n'\le n\le N$. If $n'=0$, then $E=\{x^*\}$ with
$x^*=x_0$, so $\omega=\phi_0(x_0)$ is rational with bit length polynomial
in $N$; it is the root of a linear integer polynomial with such
coefficients, and the lemma follows. Assume $n'\ge1$. Let
$R_{\mathrm{box}}\ge1$ be an integer of bit length polynomial in $N$ that
bounds the absolute values of all coordinates on $P$. Since $u^*$ is a
subvector of $x^*$, we have $\norm{u^*}_2^2\le n'R_{\mathrm{box}}^2$. Add
the ball constraint $\norm u_2^2-n'R_{\mathrm{box}}^2-1\le0$. It holds
strictly at $u^*$, keeps the optimal value $\omega$ attained at $u^*$, and
makes the feasible set compact.

We now have an objective $\psi_0(u)=\phi_0(x_0+Vu)$ and $k'\le k+1$
convex quadratic inequalities $\psi_i(u)\le0$: the substituted rows
$\phi_i(x_0+Vu)\le0$ and the ball. There are no remaining affine
constraints, and all coefficients have bitsize polynomial in $N$. Some
$\psi_i$ may be affine or constant; this does not affect the argument.
Write
\[
 \psi_i(u)=\tfrac12u^TQ_iu+a_i^Tu+c_i\qquad(0\le i\le k'),
\]
with symmetric $Q_i\succeq0$. The matrices $Q_i$ may be singular.

\emph{Step 2: regularization and elimination.}
For $0<\eta<1$, define
\begin{equation}\label{eq:regularized}
 \omega(\eta)=\min\{\psi_0(u)+\eta\norm u_2^2:
                         \psi_i(u)\le\eta\ (1\le i\le k')\}.
\end{equation}
The relaxed ball bounds every feasible point, so the feasible set is
compact, and it contains $u^*$; hence an optimum exists. The point $u^*$
satisfies every relaxed inequality strictly, so it is a Slater point of
this convex problem, and nonnegative KKT multipliers $\mu_i$ exist at
every optimum. (Strict convexity of the objective only adds uniqueness of
the optimum, which we do not need.) Stationarity reads
\begin{equation}\label{eq:matrixstationarity}
 H(\eta,\mu)\,u=-a_0-\sum_{i=1}^{k'}\mu_i a_i,\qquad
 H(\eta,\mu)=Q_0+2\eta I+\sum_{i=1}^{k'}\mu_iQ_i.
\end{equation}
The term $2\eta I$ makes $H(\eta,\mu)$ positive definite for all $\eta>0$
and all $\mu\ge0$, even when every $Q_i$ is singular. This allows us to
solve for $u$. Let
\[
 \Delta(\eta,\mu)=\det H(\eta,\mu),\qquad
 U(\eta,\mu)=-\operatorname{adj}\bigl(H(\eta,\mu)\bigr)
             \Bigl(a_0+\sum_{i=1}^{k'}\mu_ia_i\Bigr).
\]
The entries of $H(\eta,\mu)$ are affine in $(\eta,\mu)$, so $\Delta$ and
the entries of $U$ are polynomials of degree at most $n'$, and
$u=U/\Delta$. We abbreviate $\Delta=\Delta(\eta,\mu)$ and $U=U(\eta,\mu)$
and define
\begin{align}
 \Phi_i&=\tfrac12U^TQ_iU+\Delta\,a_i^TU
                  +(c_i-\eta)\Delta^2\quad(1\le i\le k'),
                  \label{eq:Gconstraint}\\
 \Phi_0&=\tfrac12U^TQ_0U+\Delta\,a_0^TU
                  +(c_0-w)\Delta^2+\eta\,U^TU.
                  \label{eq:Gobjective}
\end{align}
For $\Delta\ne0$ these are
$\Phi_i=\Delta^2(\psi_i(U/\Delta)-\eta)$ and
$\Phi_0=\Delta^2(\psi_0(U/\Delta)+\eta\norm{U/\Delta}_2^2-w)$.
Consider the polynomial system
\begin{equation}\label{eq:Ksystem}
 \Sigma(\eta,w,\mu):\quad
 \Delta>0,\quad
 \mu_i\ge0,\ \Phi_i\le0,\ \mu_i\Phi_i=0\ (1\le i\le k'),
 \quad \Phi_0=0.
\end{equation}
The condition $\Delta>0$ is implied by $\eta>0$ and $\mu\ge0$; we keep it
for clarity. For $0<\eta<1$, the system $\Sigma(\eta,w,\mu)$ has a
solution $\mu$ exactly when $w=\omega(\eta)$. If $w=\omega(\eta)$, an
optimum and its KKT multipliers give a solution: stationarity gives
$u=U/\Delta$, and feasibility, complementarity and the optimal value give
the remaining conditions. Conversely, if $\mu$ solves the system, then
$u=U/\Delta$ satisfies stationarity, feasibility and complementarity, so
$(u,\mu)$ is a KKT pair of a convex problem; hence $u$ is optimal and
$\Phi_0=0$ gives $w=\omega(\eta)$.

Every polynomial in \eqref{eq:Ksystem} has degree at most $2n'+2$. Its
coefficients have bitsize polynomial in $N$. To see this, clear the
denominators of the entries of the matrices $Q_i$ by one positive common
denominator. Each coefficient of the determinant, as a polynomial in
$(\eta,\mu)$, is then a sum of at most $n'!\,(k'+2)^{n'}$ products of at
most $n'$ integers of polynomial bitsize; since $n'\le N$ and $k'\le N+1$,
its bitsize is polynomial in $N$. The adjugate and the polynomials
$\Phi_i$ have the same property. This bounds coefficient heights even when
the number of monomials is large. Multiplying each polynomial by a
positive integer that clears its remaining denominators does not change
the sign conditions.

\emph{Step 3: the limit and its algebraic height.}
We have $\omega(\eta)\to\omega$ as $\eta\downarrow0$. For the upper
limit, $u^*$ is feasible for \eqref{eq:regularized}, so
$\omega(\eta)\le\omega+\eta\norm{u^*}_2^2$. For the lower limit, take any
sequence $\eta_j\downarrow0$ and regularized optima $u_j$. They lie in a
common ball, so a subsequence converges to some $\bar u$. The limit
satisfies $\psi_i(\bar u)\le0$ for all $i$, so
$\psi_0(\bar u)\ge\omega$, and $\omega(\eta_j)\ge\psi_0(u_j)$ gives
$\liminf_j\omega(\eta_j)\ge\omega$ along the subsequence. Applying this to
a subsequence that realizes the $\liminf$ proves the claim. No bound on
the regularized multipliers is needed.

Consequently the following formula, with one free variable $\alpha$,
defines exactly the set $\{\omega\}$:
\begin{equation}\label{eq:limitformula}
 \forall\xi\ \exists(\eta,w,\mu)\ \left[\ \xi\le0\ \lor\
   \left(
   \begin{array}{l}
    0<\eta<1,\quad\eta<\xi,\\
    -\xi<w-\alpha<\xi,\quad\Sigma(\eta,w,\mu)
   \end{array}\right)\right].
\end{equation}
Since $\xi\le0$ does not involve $\eta$, $w$ or $\mu$, this prenex
formula is equivalent to
$\forall\xi\,[\xi\le0\lor\exists(\eta,w,\mu)(\ldots)]$. For
$\alpha=\omega$, convergence gives, for every $\xi>0$, some
$0<\eta<\min\{1,\xi\}$ with $|\omega(\eta)-\omega|<\xi$, and
$w=\omega(\eta)$ with KKT multipliers $\mu$ is a witness. Conversely, if
the formula holds for $\alpha$, then witnesses for $\xi=1/j$ give
$\eta_j<1/j$ and $w_j=\omega(\eta_j)$ with $|w_j-\alpha|<1/j$; since
$\omega(\eta_j)\to\omega$, we get $\alpha=\omega$.

We apply effective block quantifier elimination
\citep[Theorem~14.16]{BasuPollackRoy2006}, also stated in
\citet[Theorem~2.27]{Basu2014}. For a prenex formula with quantifier
blocks of sizes $b_1$ and $b_2$, $\ell$ free variables, polynomials of
degree at most $D$, and integer coefficients of bitsize at most $\tau$, it
gives an equivalent quantifier-free formula whose polynomials have degree
at most $D^{O(b_1)O(b_2)}$ and integer coefficients of bitsize at most
$\tau D^{O(b_1)O(b_2)O(\ell)}$. We use this bound on coefficient heights,
not only the bound on the number of arithmetic operations. For
\eqref{eq:limitformula} the parameters are
\[
 b_1=1\ (\text{the variable }\xi),\qquad
 b_2=k'+2\le k+3\ (\text{the variables }\eta,w,\mu),\qquad
 \ell=1\ (\text{the variable }\alpha),
\]
\[
 D=2n'+2=O(n'+1),\qquad \tau=N^{O(1)}.
\]
Since $n'\le N$, we have $D^{O(k+1)}=N^{O(k+1)}$. The quantifier-free
formula therefore involves univariate integer polynomials in $\alpha$ of
degree and coefficient bitsize $N^{O(k+1)}$.

Remove the identically zero polynomials from this formula; their signs
are constant. At least one remaining polynomial vanishes at $\omega$.
Otherwise every remaining polynomial would have a constant nonzero sign on
a neighborhood of $\omega$, so the formula would have a constant truth
value there, whereas it defines the singleton $\{\omega\}$. This proves
the claim on degree and height.

Now let $\omega\ne0$. Divide that polynomial by the largest power of the
indeterminate that divides it. The quotient still vanishes at $\omega$,
has a nonzero constant term, and its coefficients are among those of the
original polynomial; let $\tau_1$ bound their bitsize. The reciprocal
polynomial has the root $1/\omega$ and a leading coefficient of absolute
value at least one, so the Cauchy bound gives
$|1/\omega|\le1+2^{\tau_1}$, that is, $|\omega|\ge1/(1+2^{\tau_1})$.
Absorbing the additive constant in the exponent proves \eqref{eq:nonzero}.
\end{proof}

The argument neither finds the active affine face algorithmically nor
needs a genericity assumption. Algebraic methods for few quadratic
constraints have direct precedents.
\citet[Theorem~1.2]{GrigorievPasechnik2005} give sampling and coefficient
bounds for zero sets of polynomials in quadratic maps; their Theorem~1.5
states an
optimization version, whose proof is deferred to a continuation of their
paper. \citet[Sections~6.2--7.2 and~8.5]{KammingaRudolph2025} give
reduced-variable formulas for optimization under a boundedness hypothesis
on the solution set, and an application to QCQP with a fixed total number
of constraints. After our affine-face reduction, squared slack variables
would turn the inequalities $\psi_i\le0$ into a bounded zero set of a
quadratic map (bounded because $u$ lies in the ball and each squared
slack equals $-\psi_i(u)$, which is bounded there), which connects the lemma to these results. We use
instead a direct convex KKT elimination and the explicit two-block limit
formula \eqref{eq:limitformula}. The affine-face reduction is what allows
arbitrarily many affine rows. The value bound itself follows from
established algebraic tools; the new part is its application to exact
penalties, with the explicit dependence on the number of nonlinear rows.

\subsection{Separation and multiplier bounds}
\begin{proof}[Proof of Theorem~\ref{thm:fixedcount}]
We bound the separation of every equality-infeasible slice from below and
a linking multiplier of every equality-feasible slice from above, each by
applying Lemma~\ref{lem:nonzerovalue} to an auxiliary problem, and then
apply Lemma~\ref{lem:slices}. Use the objective bound $|f|\le F_0$ from
Section~\ref{sec:upper}. Every auxiliary problem below is obtained by
substituting a boxed integer assignment $z$. Its coefficients keep
bitsize $O(N)$, as in \eqref{eq:height}, so its explicit input length is
at most $N^{c}$ for an absolute constant $c$, uniformly in $z$. Hence the bound $2^{-(N^{c})^{O(k+1)}}$ of
Lemma~\ref{lem:nonzerovalue} is again $2^{-N^{O(k+1)}}$.

\emph{Equality-infeasible slices.}
Let $X_z$ be a nonempty equality-infeasible slice. Let $R_0$ be a rational
bound on all residual magnitudes $|r_j(x,z)|$ over the input box, with bit
length polynomial in $N$. Minimize $t$ over the points $(x,t)$ with
$x\in X_z$ and
\[
 -t\le r_j(x,z)\le t\quad(1\le j\le m),\qquad 0\le t\le R_0.
\]
This is a nonempty compact convex quadratic problem whose only nonlinear
rows are the at most $k$ native rows that are nonlinear in $x$. Its
optimal value is $\delta_z=\min_{X_z}\norm r_\infty$, which is positive
by compactness. Lemma~\ref{lem:nonzerovalue} gives
\begin{equation}\label{eq:fixedseparation}
 \delta_z\ge2^{-N^{O(k+1)}}.
\end{equation}

\emph{Equality-feasible slices: the Slater margin.}
Let $X_z$ be an equality-feasible slice with at least one row nonlinear
in $x$. Let $P_z$ be the polyhedron defined by the native constraints that
are affine in $x$ (including the bounds) and by the relaxed equalities
$r(x,z)=0$. Define the common Slater margin of the nonlinear rows,
\begin{equation}\label{eq:slatermargin}
 \gamma_z=\max\{t:x\in P_z,\ g_i(x,z)+t\le0
           \text{ for all rows $i$ nonlinear in $x$},\ 0\le t\le1\}.
\end{equation}
The refined Slater condition provides a point of $P_z$ at which every
nonlinear row is strict, so $\gamma_z>0$. Minimizing $-t$ in this compact
convex problem with at most $k$ nonlinear rows and applying
Lemma~\ref{lem:nonzerovalue} gives
\begin{equation}\label{eq:marginbound}
 \gamma_z\ge2^{-N^{O(k+1)}}.
\end{equation}

\emph{Multipliers of the nonlinear rows.}
Let $(\bar x,\gamma_z)$ be optimal in \eqref{eq:slatermargin}, and let $x^*$ be an optimum of
the slice problem $v_z=\min\{f(x,z):x\in X_z,\ r(x,z)=0\}$. By the refined
Slater condition, $x^*$ has a KKT certificate, as in
Section~\ref{sec:upper}, with multipliers $\mu_i\ge0$ for
the nonlinear rows, nonnegative multipliers for the native rows affine in
$x$, and linking multipliers for $r=0$. The corresponding convex
Lagrangian is minimized at $x^*$, where its value is $v_z$. Evaluate it
at $\bar x$: the linking terms vanish because $r(\bar x,z)=0$, the terms
of the affine native rows are nonpositive because $\bar x\in P_z$, and
each nonlinear row contributes at most $-\gamma_z\mu_i$. Hence
\begin{equation}\label{eq:nonlinearmult}
 v_z\le f(\bar x,z)-\gamma_z\sum_i\mu_i,\qquad
 \sum_i\mu_i\le\frac{2F_0}{\gamma_z}.
\end{equation}
This is the usual bound on multipliers from a Slater point with a margin.

\emph{Linking multipliers by Cramer's rule.}
It remains to bound the linking multipliers; we do not need to bound the
multipliers of the affine native rows in this particular certificate.
The vector
\[
 b_*=-\nabla_xf(x^*,z)-\sum_i\mu_i\nabla_xg_i(x^*,z)
\]
has coordinates bounded by $2^{N^{O(1)}}(1+2F_0/\gamma_z)$, since all
gradients are affine with coefficients of polynomial bitsize and $x^*$
and $z$ lie in their boxes. By \eqref{eq:marginbound}, this bound is
$2^{N^{O(k+1)}}$. By stationarity, $b_*$ lies in the cone generated by
the gradients of the active native rows affine in $x$ and by both signs
of the linking normals $A_j^T$. By the conic Carath\'eodory theorem,
$b_*$ is a nonnegative combination of linearly independent generators
from this set; let $\nu\le n$ be their number. Their $n\times\nu$
column matrix has rank $\nu$, so some $\nu$ of its rows form a nonsingular
square submatrix. Its entries are rational input coefficients after the integer
substitution, so determinant bounds and Cramer's rule bound the entries
of its inverse by $2^{N^{O(1)}}$. Applying the inverse to the
corresponding $\nu$ coordinates of $b_*$ bounds all selected coefficients
by $2^{N^{O(k+1)}}$. This bound holds for real coefficients; the primal
point $x^*$ and the multipliers $\mu_i$ need not be rational.

Combining the coefficients of oppositely signed linking normals gives a
new KKT certificate at $x^*$: it keeps the multipliers $\mu_i$, puts the
selected coefficients on the active affine rows, and has a linking
multiplier $\lambda_z$ with
\begin{equation}\label{eq:fixedmult}
 \norm{\lambda_z}_1\le2^{N^{O(k+1)}}.
\end{equation}
Only active affine rows received multipliers and the multipliers $\mu_i$
were kept, so stationarity and complementarity still hold. As in
Section~\ref{sec:upper}, the new Lagrangian is convex and minimized at
$x^*$ with value $v_z$; dropping its nonpositive native terms on $X_z$
gives $f(x,z)+\lambda_z^Tr(x,z)\ge v_z$ for all $x\in X_z$. If the slice
has no row nonlinear in $x$, the same argument applies directly with
$b_*=-\nabla_xf(x^*,z)$; KKT conditions hold without a constraint
qualification for affine constraints, and no margin problem is needed.

\emph{Assembly.}
All constants are uniform over the integer assignments. Choose an integer
$\rho$ strictly larger than the bound in \eqref{eq:fixedmult} and than
$2F_0$ times the reciprocal of the bound in \eqref{eq:fixedseparation}.
Its bit length is $N^{O(k+1)}$. The dual norm of $\norm\cdot_\infty$ is
$\norm\cdot_1$, so Lemma~\ref{lem:slices}, with objective bounds
$\underline f=-F_0$ and $\overline f=F_0$, shows that $\rho$ is
minimizer-exact at multiplier zero for the infinity norm. The same $\rho$
works for the one-norm, because $\norm r_1\ge\norm r_\infty$ and both
norms vanish on the original feasible set.
\end{proof}

\begin{remark}
By Theorem~\ref{thm:fixedcount}, superpolynomial penalty encoding within
this model needs a growing number of nonlinear inequalities; arbitrarily
many affine rows alone do not produce it. Some growth with $k$ is
unavoidable: the chain of Remark~\ref{rem:seed} has $k=d-1$ and needs
$\Omega(N2^k)$ bits, against the $N^{O(k+1)}$ bits of
Theorem~\ref{thm:fixedcount}.
\end{remark}

\subsection{A conservative coefficient can be printed in polynomial time}
Because the encoding bounds of Theorems~\ref{thm:generalupper}
and~\ref{thm:fixedcount} are uniform, a sufficient coefficient can be
output without finding an optimum, a Slater point or an active face.

\begin{corollary}\label{cor:printing}
For each fixed continuous dimension $n$, and separately for each fixed
$k\ge0$, there is a polynomial-time algorithm with the following property.
On every input that is promised to satisfy Assumption~\ref{ass:model}
(and, in the second case, to have at most $k$ native inequalities
nonlinear in the continuous variables), it outputs in binary an integer
$\rho$ that is minimizer-exact at multiplier zero for both the infinity
norm and the one-norm.
\end{corollary}
\begin{proof}
Choose a universal integer constant $C$ such that the bounds in
Theorems~\ref{thm:generalupper} and~\ref{thm:fixedcount} are at most
\[
 B_{\mathrm{dim}}=(N+1)2^{C(n+1)}\qquad\text{and}\qquad
 B_{\mathrm{cnt}}=N^{C(k+1)},
\]
respectively. In each case some integer that is minimizer-exact at
multiplier zero has at most $B_{\mathrm{dim}}$, respectively
$B_{\mathrm{cnt}}$, bits and is therefore smaller than
$2^{B_{\mathrm{dim}}}$, respectively $2^{B_{\mathrm{cnt}}}$. By the
consequences of \eqref{eq:fixed}, every larger coefficient is
minimizer-exact at multiplier zero as well. The algorithm therefore prints
a one followed by $B_{\mathrm{dim}}$, respectively $B_{\mathrm{cnt}}$,
zeros. For fixed $n$,
respectively fixed $k$, this output and the time to produce it are
polynomial in $N$.
\end{proof}
For native constraints linear in all variables, a special case of
$k=0$, the analogous results are known: polynomial encoding length for a
jointly convex quadratic objective \citep[Theorem~11]{GuAhmedDey2020},
and polynomial-time output for MILPs with integer data
\citep[Theorem~15]{LefebvreSchmidt2025}, whose Table~1 remarks that the
extension to convex MIQPs is easy (Section~\ref{sec:related}). The case
$k=0$ of Corollary~\ref{cor:printing} also covers rows that are affine in
$x$ but contain products with integer variables, such as
$x_1z_1+z_1^2\le3$. Because all variables are bounded, such rows can
also be rewritten exactly as linear rows of polynomial size, through a
binary expansion of $z$ and exact linearization of the resulting
products, after which \citet[Theorem~11]{GuAhmedDey2020} applies when the
objective is jointly convex; this part of the corollary is therefore
mainly a convenience. The
output of Corollary~\ref{cor:printing} is a
conservative and numerically huge coefficient: the algorithm does not
verify the promises, solve the problem or approximate the least penalty,
and the constant $C$ exists by the cited effective bounds but has not
been evaluated.
